\phantomsection
\subsection*{M2C. Berry Mead}
\addcontentsline{toc}{subsection}{M2C. Berry Mead}

\textit{Berry Mead é uma categoria de inscrição para meloméis elaborados com um ou mais tipos de bagas, como framboesa, mirtilo, amora-preta, groselha (preta, vermelha e branca), morango, boysenberry, baga de sabugueiro, marionberry, amora-branca, lingonberry, huckleberry, cranberry, entre outras. De forma geral, qualquer fruta com "berry" no (no Brasil a famosa mistura de "Frutas Vermelhas": normalmente composta por morango, amora, mirtilo e framboesa) no nome se enquadra na categoria. Bagas, ou berries, podem ter sementes, mas não possuem caroço.}

\textbf{Impressão Geral}: Como um M1 Traditional Mead com um caráter frutado perceptível. Uma ampla gama de resultados é possível, dependendo das escolhas do hidromel base e dos ingredientes adicionais, mas o produto final deve ser equilibrado, agradável e prazeroso. O caráter da fruta não deve dominar completamente o hidromel.

\textbf{Aroma}: Os aromas das frutas podem ser percebidos como de um suco doce e fresco até um vinho de fruta envelhecido ou Vinho do Porto, dependendo da força alcoólica e do dulçor. Se variedades de fruta forem declaradas, o caráter varietal deve ser perceptível. Se múltiplos tipos de fruta forem declarados, eles não precisam ser percebidos em proporções iguais, tendo em vista que algumas frutas são mais intensas e distintivas que outras e podem facilmente ganhar destaque em uma mistura.

O caráter de mel pode variar de sutil a intenso, dependendo da força alcoólica e dulçor e pode expressar características do néctar de flores, refletindo quaisquer variedades de mel declaradas. Berry Meads costumam apresentar uma sugestão de acidez no aroma. O buquê de fermentação deve ser equilibrado. Hidroméis mais alcoólicos podem apresentar leves notas picantes oriundas do álcool. Os hidroméis carbonatados tendem a expressar mais aroma.

Os aromas provenientes da fruta devem complementar e realçar o hidromel base, não dominar.

\textbf{Aparência}: Limpidez de alta a brilhante. Quanto maior o brilho, mais desejável. A presença de bolhas, espuma, colarinho ou efervescência deve corresponder ao nível de carbonatação declarado. A cor é derivada da variedade de fruta e de mel utilizadas. Quanto mais vivas e brilhantes as cores, mais desejável.

\textbf{Sabor}: Semelhante ao Aroma quanto a característica e intensidade de fruta e mel, sendo mais proeminente em versões mais alcoólicas e doces. O caráter da fruta deve ser perceptível e em equilíbrio com o nível de dulçor do hidromel. Algumas frutas podem conferir gostos básicos (amargo, doce, ácido) além de suas características típicas. Muitas frutas conferem acidez e taninos naturais, que precisam estar adequadamente em equilíbrio com o hidromel base.

Dulçor residual e final de boca devem corresponder ao nível de dulçor declarado. A acidez deve ser equilibrada, macia ou vibrante, mas não agressiva. Taninos podem fazer um hidromel adocicado parecer mais seco. Características ásperas, desagradáveis ou excessivamente sulfurosas ou notas de levedura/autólise são indesejáveis.

\textbf{Sensação na Boca}: A carbonatação segue o nível declarado, variando de tranquilo a espumante. O corpo geralmente aumenta com o dulçor e a força alcoólica, tipicamente de médio-baixo a alto. A percepção de álcool acompanha a força alcoólica declarada, variando de ausente a perceptível e com sensação de aquecimento. Os taninos de frutas mais escuras podem conferir corpo e uma sensação de secura, mas não deve ser excessivo.

\textbf{Ingredientes}: Os mesmos ingredientes de um M1 Traditional Mead, com a adição de pelo menos um tipo de fruta de baga.

\textbf{Comentários}: Um Berry Mead que também contenha outra fruta deve ser inscrito como M2E Other Fruit Mead. Versões com especiarias devem ser inscritas como M3C Fruit and Spice Mead. Hidroméis com frutas e outros ingredientes que não forem frutas, quando perceptíveis, devem ser inscritos na categoria M4F (Experimental Mead).

\textbf{Instruções para Inscrição}: Participantes devem obrigatoriamente especificar os níveis de dulçor, carbonatação e força alcoólica. As variedades de mel podem ser declaradas. A fruta deve obrigatoriamente ser declarada.

\textbf{Exemplos Comerciais}: Moonlight Wild, Prairie Rose Honeyberry, Rabbit's Foot Raspberry Mead, Redstone Boysenberry Nectar, Schramm's Mead Black Heart, St. Ambrose Rhythm and Blues, White Winter Strawberry Mead.
